% !PRINTED 206
\weberpage{206}

The common denominator of the fractions occurring in (1):

\[
\prod^{\nu} \vartheta_{g_1,g_2}\left(u + \frac{2\nu}{n}\right) = \pm \prod^{\nu} \vartheta_{g_1,g_2}\left(u + \frac{\nu}{n}\right)
\]

is, by the last formulae of § 33, equal, apart from a constant
factor, to

\[
\vartheta_{g_1,g_2}(nu, n\omega),
\]

and if we therefore set the sum (1) equal to

\weq{2}{\frac{\Phi(u)}{\vartheta_{g_1,g_2}(nu, n\omega)}}

then $\Phi(u)$ is an \gesperrt{entire} transcendental function of $u$.
Now we have

\[
\frac{\vartheta_{11}\left(u + \frac{2\nu+1}{u}\right)}{\vartheta_{g_1,g_2}\left(u + \frac{2\nu+1}{u}\right)} = (-1)^{g_1+1} \frac{\vartheta_{11}\left(u + \frac{2\nu+n+1}{n}\right)}{\vartheta_{g_1,g_2}\left(u + \frac{2\nu+n+1}{n}\right)}
\]
% [sic] The German source prints u in both left-hand inner denominators; n is intended.

and $\nu$ runs, simultaneously with $\nu + \frac{1}{2}(n+1)$, through a complete system of residues with respect to $n$. From this it follows by (1) and (2) that

\weq{3}{\Phi\left(u + \frac{1}{n}\right) = -e^{-\frac{4\nu'\pi i}{n}} \Phi(u)}

and similarly

\weq{4}{\Phi(u + \omega) = -e^{-n\pi i(2u+\omega)} \Phi(u).}

From this it follows that the function $\Phi(u)$ is a $t$-function, that is, a theta-function, % [source typo] German prints t; theta is intended.
with the periods $1/n$, $\omega$, one that is determined by the conditions (3), (4)
up to a factor independent of $u$ and can be expressed by
a $\vartheta$-function (§ 19). One sees easily that
the conditions (3), (4) are satisfied by the function

\[
e^{-4\pi i\nu' u} \vartheta_{11}(nu - 4\nu'\omega, n\omega)
\]
% [sic] The German source prints 4 here but 2 in the corresponding argument in (5).

so that there results:

\weq{5}{\sum^{\nu} e^{\frac{8\nu\nu'\pi i}{n}} \frac{\vartheta_{11}\left(u + \frac{2\nu}{n}\right)}{\vartheta_{g_1,g_2}\left(u + \frac{2\nu}{n}\right)} = Ce^{-4\pi i\nu' u} \frac{\vartheta_{11}(nu - 2\nu'\omega, n\omega)}{\vartheta_{g_1,g_2}(nu, n\omega)}.}

The determination of the constant $C$ can easily be carried out if one multiplies on the right and on the left by

\[
\vartheta_{g_1,g_2}(u)
\]

and then

\[
u = \frac{(1-g_2) + (1-g_1)\omega}{2}
\]
%CONT>>
